A Figura~\ref{fig:cronograma} apresenta o cronograma entre outubro de 2026 e junho de 2027, em quinzenas. Cinza indica atividades realizadas, incluindo a ampliação GK-P2.6-dev; azul indica continuidade proposta. Datas futuras dependem do orientador e dos revisores. O número de oportunidades será definido no novo plano, e a execução final depende da ratificação das referências e do congelamento do protocolo.

Quando atividades dependentes compartilham uma quinzena, a etapa posterior somente se inicia após o encerramento formal da anterior. Isso se aplica à conferência e ratificação antes do congelamento (atividades 8 a 10), à execução antes dos pareceres (11 e 12) e ao encerramento dos pareceres e da conciliação antes do cálculo dos indicadores (12 e 13). Na atividade 6, o trecho cinza corresponde ao início realizado na primeira quinzena de novembro de 2026; o trecho azul indica a continuidade proposta até dezembro.

\begin{figure}[p]
\centering
\caption{Cronograma das atividades do projeto em intervalos quinzenais}
\label{fig:cronograma}
\begin{ganttchart}[
    y unit title=0.5cm,
    y unit chart=0.32cm,
    x unit=0.78cm,
    vgrid,hgrid,
    title height=1,
    title/.style={fill=none,draw=black!65},
    title label font=\bfseries\scriptsize,
    bar/.append style={fill=cyan!35,draw=black!65},
    bar height=0.65,
    bar label font=\scriptsize,
    progress label text={},
    inline=false
]{1}{18}
\gantttitle{2026}{6}\gantttitle{2027}{12}\\
\gantttitle{Out}{2}\gantttitle{Nov}{2}\gantttitle{Dez}{2}
\gantttitle{Jan}{2}\gantttitle{Fev}{2}\gantttitle{Mar}{2}
\gantttitle{Abr}{2}\gantttitle{Mai}{2}\gantttitle{Jun}{2}\\
\gantttitlelist{1,2}{1}\gantttitlelist{1,2}{1}\gantttitlelist{1,2}{1}
\gantttitlelist{1,2}{1}\gantttitlelist{1,2}{1}\gantttitlelist{1,2}{1}
\gantttitlelist{1,2}{1}\gantttitlelist{1,2}{1}\gantttitlelist{1,2}{1}\\
\ganttbar[bar/.append style={fill=gray!55}]{1)}{2}{2}\\
\ganttbar[bar/.append style={fill=gray!55}]{2)}{2}{3}\\
\ganttbar[bar/.append style={fill=gray!55}]{3)}{2}{3}\\
\ganttbar[bar/.append style={fill=gray!55}]{4)}{2}{2}\\
\ganttbar[bar/.append style={fill=gray!55}]{5)}{3}{3}\\
\ganttbar[bar/.append style={fill=gray!55}]{6)}{3}{3}
\ganttbar{}{4}{6}\\
\ganttbar{7)}{5}{6}\\
\ganttbar{8)}{9}{9}\\
\ganttbar{9)}{9}{10}\\
\ganttbar{10)}{10}{10}\\
\ganttbar{11)}{11}{11}\\
\ganttbar{12)}{11}{13}\\
\ganttbar{13)}{13}{14}\\
\ganttbar{14)}{14}{16}\\
\ganttbar{15)}{15}{17}\\
\ganttbar{16)}{17}{18}
\end{ganttchart}

\par\smallskip
{\footnotesize Cinza: realizado. Azul: proposto. 1 e 2: primeira e segunda quinzenas.\par}
\medskip
\begin{minipage}{\textwidth}
\small\SingleSpacing
\begin{enumerate}[label=\arabic*),nosep,leftmargin=0.8cm]
\item Consolidação do tema, do problema e dos objetivos.
\item Levantamento bibliográfico e seleção do acervo.
\item Implementação da prova de conceito.
\item Execução do piloto histórico.
\item Verificações funcionais, lotes de desenvolvimento e preparação do protocolo e dos instrumentos.
\item Redação e revisão do Projeto de Pesquisa/TCC I.
\item Preparação e apresentação da proposta de TCC I.
\item Especificação das duas proteções, adaptação do experimento e revisão dos casos para as dez matérias.
\item Confirmação dos avaliadores e ratificação humana das fontes, dos casos e da rubrica.
\item Congelamento do protocolo, código, índice e modelos.
\item Execução das oportunidades do novo plano final, após congelamento.
\item Pareceres humanos e conciliação das divergências.
\item Cálculo dos indicadores de proteção e bloqueio indevido; análise complementar dos estilos.
\item Análise, discussão dos resultados e conclusão.
\item Revisão final do texto e das referências.
\item Preparação da apresentação e entrega final.
\end{enumerate}
\end{minipage}
\fonte{Elaboração própria; atividades futuras propostas, sem prazo institucional presumido.}
\end{figure}
